\documentclass[12pt,oneside]{book}

%==================================================================
% METADATOS
%==================================================================
\newcommand{\universidad}{Universidad de Buenos Aires (UBA)}
\newcommand{\carrera}{Carrera de Derecho}
\newcommand{\materia}{Derecho Constitucional Profundizado}
\newcommand{\titulo}{Apuntes de Derecho Constitucional Profundizado}
\newcommand{\subtitulo}{Segunda parte: repaso y cierre previo al segundo parcial}
\newcommand{\autor}{Isaac Edgar Camacho Ocampo}
\newcommand{\anio}{2026}

%==================================================================
% PAQUETES
%==================================================================
\usepackage[utf8]{inputenc}
\usepackage[T1]{fontenc}
\usepackage[spanish,es-nodecimaldot]{babel}
\usepackage[
    a4paper,
    top=2.5cm,
    bottom=2.5cm,
    left=3cm,
    right=2.5cm,
    headheight=15pt
]{geometry}
\usepackage{amsmath}
\usepackage{amssymb}
\usepackage{mathtools}
\usepackage{graphicx}
\usepackage{float}
\usepackage{caption}
\usepackage{subcaption}
\usepackage{booktabs}
\usepackage{array}
\usepackage{longtable}
\usepackage{tabularx}
\usepackage{enumitem}
\usepackage{xcolor}
\usepackage[most]{tcolorbox}
\usepackage{fancyhdr}
\usepackage{ragged2e}
\usepackage[
    colorlinks=true,
    linkcolor=blue!50!black,
    citecolor=green!40!black,
    urlcolor=blue!60!black,
    pdftitle={\titulo},
    pdfauthor={\autor},
    pdfsubject={Apuntes universitarios},
    bookmarks=true
]{hyperref}

%==================================================================
% CONFIGURACIÓN
%==================================================================
\graphicspath{
    {./img/}
    {./graficos/}
    {./figuras/}
}
\setcounter{tocdepth}{2}
\setlength{\parindent}{0pt}
\setlength{\parskip}{0.5em}
\setlist[itemize]{
    topsep=0.3em,
    itemsep=0.2em
}
\setlist[enumerate]{
    topsep=0.3em,
    itemsep=0.2em
}
\clubpenalty=10000
\widowpenalty=10000

%==================================================================
% COMANDOS
%==================================================================
\newcommand{\abs}[1]{\left|#1\right|}
\newcommand{\norm}[1]{\left\lVert#1\right\rVert}
\newcommand{\vect}[1]{\mathbf{#1}}
\newcommand{\deriv}[2]{\frac{d#1}{d#2}}
\newcommand{\pderiv}[2]{\frac{\partial#1}{\partial#2}}

%==================================================================
% ENTORNOS / CAJAS
%==================================================================
\newtcolorbox{definition}[1]{
    enhanced,
    breakable,
    title={Definición: #1},
    fonttitle=\bfseries,
    colback=blue!3,
    colframe=blue!50!black,
    boxrule=0.8pt,
    arc=2mm
}

\newtcolorbox{important}{
    enhanced,
    breakable,
    title={Importante},
    fonttitle=\bfseries,
    colback=yellow!8,
    colframe=orange!70!black,
    boxrule=0.8pt,
    arc=2mm
}

\newtcolorbox{example}[1]{
    enhanced,
    breakable,
    title={Ejemplo: #1},
    fonttitle=\bfseries,
    colback=green!4,
    colframe=green!40!black,
    boxrule=0.8pt,
    arc=2mm
}

\newtcolorbox{formula}[1]{
    enhanced,
    breakable,
    title={Fórmula / Regla: #1},
    fonttitle=\bfseries,
    colback=purple!4,
    colframe=purple!50!black,
    boxrule=0.8pt,
    arc=2mm
}

\newtcolorbox{exercise}[1]{
    enhanced,
    breakable,
    title={Ejercicio: #1},
    fonttitle=\bfseries,
    colback=gray!5,
    colframe=gray!60!black,
    boxrule=0.8pt,
    arc=2mm
}

\newtcolorbox{note}{
    enhanced,
    breakable,
    title={Nota},
    fonttitle=\bfseries,
    colback=white,
    colframe=gray!50,
    boxrule=0.6pt,
    arc=2mm
}

%==================================================================
% CONFIGURACIÓN FINAL
%==================================================================
\pagestyle{fancy}
\fancyhf{}
\fancyhead[L]{\nouppercase{\leftmark}}
\fancyhead[R]{\materia}
\fancyfoot[C]{\thepage}
\renewcommand{\headrulewidth}{0.4pt}
\renewcommand{\footrulewidth}{0pt}

\fancypagestyle{plain}{
    \fancyhf{}
    \fancyfoot[C]{\thepage}
    \renewcommand{\headrulewidth}{0pt}
}

%==================================================================
\begin{document}
%==================================================================

%------------------------------------------------------------------
% PORTADA
%------------------------------------------------------------------
\begin{titlepage}

\thispagestyle{empty}

\begin{center}

\vspace*{1cm}

{\large \textsc{\universidad}}

\vspace{0.2cm}

{\large \carrera}

\vspace{3cm}

{\Huge \textbf{\materia}}

\vspace{1cm}

{\LARGE \subtitulo}

\vspace{0.8cm}

\textit{Estudiar con comprensión es construir conocimiento duradero.}

\vspace{1.2cm}

\rule{0.70\textwidth}{0.5pt}

\vspace{0.5cm}

{\large Apuntes de estudio}

\vspace{0.5cm}

\rule{0.70\textwidth}{0.5pt}

\vfill

{\large \autor}

\vspace{0.3cm}

{\large \anio}

\end{center}

\vfill

\begin{flushleft}
\textit{Este es un modesto aporte a los estudiantes de ``\materia'' no pretende sustituir el material oficial de la materia.}
\end{flushleft}

\end{titlepage}

%------------------------------------------------------------------
% ÍNDICE
%------------------------------------------------------------------
\tableofcontents
\clearpage

%==================================================================
% CONTENIDO
%==================================================================

%------------------------------------------------------------------
\chapter{Libertad de culto, réplica y prensa}
%------------------------------------------------------------------

\section{Libertad religiosa y sostenimiento del culto}

El estándar constitucional compatibiliza la \textbf{libertad religiosa}, es decir, el derecho a profesar los cultos reconocidos dentro del país, con una especial protección de la \textbf{no discriminación}: cualquier división o categorización jurídica basada en las religiones que se profesan queda comprendida dentro de las denominadas \textbf{categorías sospechosas}.

La libertad religiosa es un principio cardinal de la Constitución. Sin embargo, se mantiene el sostenimiento del culto católico apostólico romano, que tiene un \textbf{carácter eminentemente económico}.

\begin{quote}
\textbf{Constitución Nacional, art. 2.} El Gobierno federal sostiene el culto católico apostólico romano.
\end{quote}

Este sostenimiento tiene un origen histórico, que se plasma también en el reconocimiento específico de la Iglesia como una categoría de \textbf{persona jurídica pública} con algunas prerrogativas, como la \textbf{inembargabilidad de los bienes afectados al culto activo}.

\subsection{El fallo Castillo}

El fallo \textbf{Castillo} es una sentencia actual de la Corte que ofrece una aproximación práctica a la manera de entender la libertad religiosa y, al mismo tiempo, de admitir el sostenimiento del culto católico. Se refiere a la educación en la provincia de Salta.

Los fundamentos de la Corte pueden ordenarse del siguiente modo:
\begin{itemize}
    \item Si bien existe un sostenimiento del culto católico, ello \textbf{no implica un privilegio ni un trato inequitativo} por el cual la religión católica deba tener una especial predilección en las políticas públicas, sobre todo en materia educativa.
    \item Más allá de los orígenes históricos que lo sustentan, el sostenimiento del culto es de carácter eminentemente económico, en cuanto a la asistencia del gobierno y a una categoría específica que tiene la Iglesia. En la práctica de la educación, y además de la educación inicial, no puede suponer una prevalencia de la religión católica.
    \item La ley provincial establecía formalmente que se impartía la materia de religión mediante un muestreo de las distintas religiones practicadas en el país, esto es, una religión básica de todas las religiones con presencia en el país.
    \item La prueba en la causa demostró que, por cuestiones de medios, de capacitación de profesores y de material de enseñanza, solo se impartía catequesis de la religión católica apostólica romana.
    \item Ello supone una suerte de \textbf{monopolio} de la religión católica en la educación referida a temas religiosos, y un trato inequitativo contrario al régimen constitucional de libertad de culto.
    \item El sistema, tal como se aplicaba en Salta, fue declarado \textbf{inconstitucional}. Las políticas públicas no pueden privilegiar una religión por sobre otra.
\end{itemize}

\begin{important}
\textbf{Sostener el culto no es privilegiarlo.} El sostenimiento del culto católico tiene carácter eminentemente económico; no autoriza una prevalencia de la religión católica en las políticas públicas, sobre todo en la educación.
\end{important}

\section{Derecho de réplica}

El \textbf{derecho de réplica} es un mecanismo complementario de la libertad de expresión, que permite, en igualdad de condiciones, dar voz a las posturas disidentes o, muchas veces, minoritarias.

Su recepción en el derecho argentino se produce a través del \textbf{Pacto de San José de Costa Rica}, norma internacional que internaliza este derecho.

\begin{quote}
\textbf{Convención Americana sobre Derechos Humanos (Pacto de San José de Costa Rica), art. 14.}
\begin{enumerate}
    \item Toda persona afectada por informaciones inexactas o agraviantes emitidas en su perjuicio a través de medios de difusión legalmente reglamentados y que se dirijan al público en general, tiene derecho a efectuar por el mismo órgano de difusión su rectificación o respuesta en las condiciones que establezca la ley.
    \item En ningún caso la rectificación o la respuesta eximirán de las otras responsabilidades legales en que se hubiese incurrido.
    \item Para la efectiva protección de la honra y la reputación, toda publicación o empresa periodística, cinematográfica, de radio o televisión tendrá una persona responsable que no esté protegida por inmunidades ni disponga de fuero especial.
\end{enumerate}
\end{quote}

En el fallo \textbf{Ekmekdjian} la Corte vinculó este derecho con el concepto de \textbf{operatividad de los tratados internacionales}: no se necesita una ley interna que reglamente el pacto interamericano, porque este es suficientemente operativo, y ninguna interpretación constitucional podría menoscabar su operatividad.

\section{Libertad de prensa}

La libertad de prensa es el complemento del derecho de réplica y tiene dos vertientes:

\begin{table}[H]
\centering
\renewcommand{\arraystretch}{1.3}
\begin{tabularx}{\textwidth}{
    >{\raggedright\arraybackslash}p{0.17\textwidth}
    >{\raggedright\arraybackslash}p{0.27\textwidth}
    >{\raggedright\arraybackslash}X
}
\toprule
\textbf{Dimensión} & \textbf{Titular} & \textbf{Qué protege} \\
\midrule
\textbf{Social activa} & Todos los hombres y mujeres que emiten información u opinión & Expresarse sin censura previa, sobre opiniones personales y políticas, en los distintos medios. \\
\textbf{Pasiva} & Quien recibe la información & Percibir distintos tipos de informaciones, escuchando voces disidentes, con un panorama amplio y sin voces silenciadas ni opacadas. \\
\bottomrule
\end{tabularx}
\caption{Las dos dimensiones de la libertad de prensa}
\end{table}

La libertad de prensa es \textbf{piedra angular del sistema republicano, representativo y democrático}: al asegurar en un foro público la pluralidad de voces, asegura implícitamente el contralor de los actos de gobierno. Implica también asegurar el acceso a posiciones disidentes de la gestión pública predominante, lo que redunda en la \textbf{alternancia} en la gestión política (referencia al fallo de Formosa). Estas dos caras de la misma moneda ayudan a operativizar esa alternancia: el derecho de todos a expresar voces disidentes y el derecho de todos a recibir esa información y acceder a ella.

\subsection{Prohibición casi absoluta de la censura previa}

Otra piedra angular de la libertad de prensa es la \textbf{prohibición casi absoluta de la censura previa}. Solo cuando están comprometidos intereses o derechos de menores de edad o de sectores muy vulnerables podría operar, con muchas salvaguardas y siempre en presupuestos muy específicos, alguna prohibición de publicación. La regla general es que la reparación es \textbf{ex post}, consecuencia de una eventual conducta ilegítima en el ejercicio de la libertad de prensa.

\begin{important}
\textbf{Reparación ex post, no censura previa.} En materia de libertad de prensa se trabaja siempre bajo un criterio de reparación posterior, sin lugar para decisiones de censura previa.
\end{important}

\begin{note}
Los temas \textbf{Real Malicia} (derivada del fallo New York Times v. Sullivan) y \textbf{doctrina Campillay} corresponden a la primera parte de la materia y ya fueron evaluados; en esta segunda parte la libertad de prensa se toma a nivel conceptual, como complemento del derecho de réplica.
\end{note}

%------------------------------------------------------------------
\section{Cuadro Cornell: libertad de culto, réplica y prensa}
%------------------------------------------------------------------

\begin{longtable}{@{}p{0.26\textwidth}p{0.70\textwidth}@{}}
\toprule
\textbf{Preguntas / palabras clave} & \textbf{Notas / respuestas} \\
\midrule
\endfirsthead
\toprule
\textbf{Preguntas / palabras clave} & \textbf{Notas / respuestas} \\
\midrule
\endhead

\textbf{¿Cómo se compatibiliza la libertad religiosa con el sostenimiento del culto católico?} &
La libertad religiosa es un principio cardinal: protege el derecho a profesar los cultos reconocidos y exige especial protección contra la discriminación (toda categorización jurídica basada en la religión es una categoría sospechosa). El sostenimiento del culto católico (art. 2) se mantiene con carácter eminentemente económico y origen histórico, y la Iglesia es persona jurídica pública con prerrogativas como la inembargabilidad de los bienes afectados al culto activo. \\[0.8em]

\textbf{¿Qué resolvió la Corte en el fallo Castillo?} &
Declaró inconstitucional el sistema de enseñanza religiosa de Salta: la ley decía impartir una religión básica de todas las religiones presentes en el país, pero la prueba demostró que solo se enseñaba catequesis católica. Eso es un monopolio y un trato inequitativo; las políticas públicas no pueden privilegiar una religión sobre otra. \\[0.8em]

\textbf{¿Qué es el derecho de réplica y por qué es operativo?} &
Mecanismo complementario de la libertad de expresión que, en igualdad de condiciones, da voz a posturas disidentes o minoritarias. Lo recibe el Pacto de San José de Costa Rica (art. 14). Según el fallo Ekmekdjian, no hace falta una ley interna que lo reglamente: es suficientemente operativo y ninguna interpretación constitucional puede menoscabar esa operatividad. \\[0.8em]

\textbf{¿Qué dimensiones tiene la libertad de prensa y qué papel cumple la censura previa?} &
Dimensión social activa (todos pueden expresar opiniones sin censura previa) y dimensión pasiva (todos pueden recibir información plural). Es piedra angular del sistema republicano: asegura el contralor de los actos de gobierno y la alternancia. La censura previa está casi absolutamente prohibida; solo con muchas salvaguardas podría operar para proteger a menores o a sectores muy vulnerables. La regla es la reparación ex post. \\

\midrule
\multicolumn{2}{@{}p{0.96\textwidth}@{}}{\textbf{Resumen:} La libertad religiosa convive con el sostenimiento económico del culto católico siempre que no haya privilegio ni discriminación (Castillo). El derecho de réplica, plenamente operativo por el Pacto de San José (Ekmekdjian), complementa una libertad de prensa de doble dimensión, cuya regla es la reparación posterior y la prohibición casi absoluta de la censura previa.} \\
\bottomrule
\end{longtable}

%------------------------------------------------------------------
\chapter{Propiedad y expropiación}
%------------------------------------------------------------------

\section{Derecho de propiedad}

Según la jurisprudencia de la Corte, nadie tiene derecho al mantenimiento de leyes o reglamentos; pero sí tiene derecho a una \textbf{reparación} si, por las modificaciones de las normas a través del tiempo, se modifica un derecho adquirido al amparo de la norma anterior.

El \textbf{artículo 17} consagra el derecho de propiedad entendido como cualquier derecho o interés conexo con la propia vida y el cuerpo de la persona, y tutela cualquier tipo de derechos o intereses jurídicos, por ejemplo:
\begin{itemize}
    \item propiedades inmuebles y muebles;
    \item derechos adquiridos;
    \item derechos de propiedad intelectual, marcas y patentes, es decir, derechos de titularidad por la creación de algo;
    \item el respeto a los contratos suscriptos por acuerdo de voluntades legítimas.
\end{itemize}

\begin{quote}
\textbf{Constitución Nacional, art. 17.} La propiedad es inviolable, y ningún habitante de la Nación puede ser privado de ella, sino en virtud de sentencia fundada en ley. La expropiación por causa de utilidad pública, debe ser calificada por ley y previamente indemnizada. Sólo el Congreso impone las contribuciones que se expresan en el artículo 4º. Ningún servicio personal es exigible, sino en virtud de ley o de sentencia fundada en ley. Todo autor o inventor es propietario exclusivo de su obra, invento o descubrimiento, por el término que le acuerde la ley. La confiscación de bienes queda borrada para siempre del Código Penal argentino. Ningún cuerpo armado puede hacer requisiciones, ni exigir auxilios de ninguna especie.
\end{quote}

De este artículo resulta que el desapoderamiento legítimo de un bien debe estar precedido de una \textbf{indemnización previa} y del cumplimiento del \textbf{debido proceso}. Nadie puede ser despojado de su propiedad \emph{manu militari}: debe existir una declaración de utilidad pública, una ley declarativa y una determinación que lleve luego al Poder Ejecutivo a iniciar la expropiación.

\subsection{El fallo Bourdieu}

\textbf{Bourdieu} es el fallo clásico donde la Corte sienta el concepto de derecho de propiedad. Está citado en el programa y, según lo indicado en clase, es un fallo de más de cien años, de la década de 1920.
% TODO: contenido incompleto o ambiguo en las notas originales (fecha aproximada del fallo Bourdieu: <<década del 20, del 25, si no me equivoco>>)

\section{Ley de expropiación}

Se deben comprender los conceptos generales de la ley de expropiación.

\begin{itemize}
    \item La ley de declaración de utilidad pública debe \textbf{individualizar} el bien o los bienes objeto de expropiación, o dar las pautas para que se individualicen.
    \item El \textbf{pago previo} es una novedad, entre comillas, de la Constitución, que lo prevé taxativamente.
\end{itemize}

\subsection{Expropiación inversa}

\begin{definition}{Expropiación inversa}
Es aquella en la que no es el Estado quien pide al juez que se expropie y paga por el bien que desapodera al particular, sino que es el \textbf{particular} quien la solicita porque, producto de alguna regulación pública, no puede hacer un uso y goce racional y habitual de su propiedad. Siempre es el particular el que pide la expropiación.
\end{definition}

\begin{example}{La estación de trenes frente a la casa}
Una persona tiene una casa y frente a ella se instala una estación de trenes, con una entrada y salida muy dificultosa a su propiedad. Si el particular demuestra que ya no puede usar y gozar de su propiedad normalmente, podría llegar a darse la expropiación inversa.
\end{example}

\begin{table}[H]
\centering
\renewcommand{\arraystretch}{1.3}
\begin{tabularx}{\textwidth}{
    >{\raggedright\arraybackslash}p{0.18\textwidth}
    >{\raggedright\arraybackslash}X
    >{\raggedright\arraybackslash}X
}
\toprule
\textbf{Aspecto} & \textbf{Expropiación} & \textbf{Expropiación inversa} \\
\midrule
\textbf{Quién la pide} & El Estado, ante el juez & El particular afectado \\
\textbf{Qué ocurre} & El Estado desapodera el bien y paga por él & Una regulación pública impide al particular hacer un uso y goce racional y habitual de su propiedad \\
\textbf{Ejemplo} & Utilidad pública declarada por ley sobre un bien individualizado & Una estación de trenes frente a la casa dificulta la entrada y salida a la propiedad \\
\bottomrule
\end{tabularx}
\caption{Expropiación y expropiación inversa}
\end{table}

\subsection{Indemnización y proyecto de ley}

Según la ley actual y la ley histórica, vigente desde hace más de cien años, solo se indemniza el \textbf{daño emergente} directo de la expropiación; no se indemniza el lucro cesante, la pérdida de chance ni ningún otro ítem referido a pérdidas futuras.

Se encuentra en discusión un proyecto de ley que, de aprobarse, ampliaría la indemnización: ya no se indemnizaría solo el daño emergente, sino que podría ser objeto de reparación también el \textbf{lucro cesante}.

\begin{table}[H]
\centering
\renewcommand{\arraystretch}{1.3}
\begin{tabularx}{\textwidth}{
    >{\raggedright\arraybackslash}p{0.25\textwidth}
    >{\raggedright\arraybackslash}X
    >{\raggedright\arraybackslash}X
}
\toprule
 & \textbf{Ley actual e histórica} & \textbf{Proyecto en discusión} \\
\midrule
\textbf{Daño emergente directo} & Se indemniza & Se indemnizaría \\
\textbf{Lucro cesante} & No se indemniza & Podría llegar a ser objeto de reparación \\
\textbf{Pérdida de chance y pérdidas futuras} & No se indemnizan & No fue mencionado en la clase \\
\bottomrule
\end{tabularx}
\caption{Qué se indemniza hoy y qué prevé el proyecto}
\end{table}

%------------------------------------------------------------------
\section{Cuadro Cornell: propiedad y expropiación}
%------------------------------------------------------------------

\begin{longtable}{@{}p{0.26\textwidth}p{0.70\textwidth}@{}}
\toprule
\textbf{Preguntas / palabras clave} & \textbf{Notas / respuestas} \\
\midrule
\endfirsthead
\toprule
\textbf{Preguntas / palabras clave} & \textbf{Notas / respuestas} \\
\midrule
\endhead

\textbf{¿Qué protege el derecho de propiedad del art. 17 y qué aporta el fallo Bourdieu?} &
Cualquier derecho o interés jurídico: inmuebles, muebles, derechos adquiridos, propiedad intelectual, marcas y patentes, contratos legítimamente celebrados. Nadie tiene derecho al mantenimiento de leyes, pero sí a reparación si una reforma modifica un derecho adquirido. Bourdieu es el fallo clásico que fija el concepto de propiedad. \\[0.8em]

\textbf{¿Cuáles son los requisitos de la expropiación?} &
Declaración de utilidad pública por ley, que individualice el bien o dé pautas para hacerlo; indemnización previa y debido proceso. No cabe el despojo \emph{manu militari}. \\[0.8em]

\textbf{¿Qué es la expropiación inversa?} &
Es la que pide el particular cuando, por una regulación pública, no puede hacer un uso y goce racional y habitual de su propiedad (ejemplo: una estación de trenes frente a la casa). \\[0.8em]

\textbf{¿Qué se indemniza en la expropiación?} &
Hoy se indemniza solo el daño emergente directo; no el lucro cesante ni la pérdida de chance. Un proyecto de ley en discusión incluiría el lucro cesante. \\

\midrule
\multicolumn{2}{@{}p{0.96\textwidth}@{}}{\textbf{Resumen:} El artículo 17 protege cualquier derecho o interés jurídico y exige ley, indemnización previa y debido proceso para expropiar. La expropiación inversa es la vía del particular afectado por regulaciones públicas. Hoy solo se indemniza el daño emergente directo.} \\
\bottomrule
\end{longtable}

%------------------------------------------------------------------
\chapter{Derechos colectivos y acciones de clase}
%------------------------------------------------------------------

\section{Legitimación activa y fallo Halabi}

La \textbf{legitimación activa} debe entenderse como un presupuesto para que exista una \textbf{cuestión justiciable}.

En el fallo \textbf{Halabi}, la Corte sentó la distinción preponderante entre tres categorías de derechos:

\begin{table}[H]
\centering
\renewcommand{\arraystretch}{1.3}
\begin{tabularx}{\textwidth}{
    >{\raggedright\arraybackslash}p{0.34\textwidth}
    >{\raggedright\arraybackslash}X
}
\toprule
\textbf{Categoría} & \textbf{Características según lo explicado en clase} \\
\midrule
\textbf{Derecho individual subjetivo} & Derecho propio de una persona determinada. \\
\textbf{Derechos de incidencia colectiva} & Derechos colectivos, que exceden lo individual. \\
\textbf{Intereses individuales homogéneos} & Tienen una causa común con otro tipo de pretensiones o reclamaciones. Este grupo da lugar a las acciones de clase. \\
\bottomrule
\end{tabularx}
\caption{Las tres categorías de derechos del fallo Halabi}
\end{table}

\begin{quote}
\textbf{Constitución Nacional, art. 43, segundo párrafo.} Podrán interponer esta acción contra cualquier forma de discriminación y en lo relativo a los derechos que protegen al ambiente, a la competencia, al usuario y al consumidor, así como a los derechos de incidencia colectiva en general, el afectado, el defensor del pueblo y las asociaciones que propendan a esos fines, registradas conforme a la ley, la que determinará los requisitos y formas de su organización.
\end{quote}

\section{Acciones de clase}

\begin{definition}{Acción de clase}
Es una acción que puede iniciar el \textbf{defensor del pueblo}, las \textbf{asociaciones} de defensa del bien jurídico de que se trate y, en última instancia, el \textbf{afectado} en representación de su clase. Versa sobre derechos que tienen una \textbf{causa homogénea}: si bien son divisibles en cuanto a la cuantificación del daño, tienen un origen común, fundado en una causa que alcanza a una pluralidad de sujetos.
\end{definition}

\begin{example}{La planta automotriz y el vicio latente}
Una planta de montaje de una automotriz que durante seis meses produce camionetas con una falla en una pieza o en un área del rodado (un vicio latente) que puede acortar el buen funcionamiento del vehículo. Ese hecho común generador de un daño plural es lo que da su fisionomía particular a la acción de clase.
\end{example}

\subsection{Utilidad de la acción de clase}

\begin{itemize}
    \item Se la ha catalogado como una herramienta útil para \textbf{mantener a raya a las grandes corporaciones}: sujetos económicos de mucho alcance y poder, que brindan un servicio masivo. Es un freno a quien tiene un mercado cuasi cautivo.
    \item Se caracteriza por reclamos que, individualmente, quizás no tengan mucho peso. Para disciplinar a las corporaciones, el abogado puede fundar una clase, representarla y obtener un monto muy atractivo si litiga contra estas asociaciones mayoritarias y de gran poder económico.
    \item Se busca que las acciones de clase se proyecten y sean moneda corriente en la impugnación de este tipo de actos.
    \item Nunca puede ser iniciada por el propio Estado para cuestionar un punto de su gestión: no podría ser cuestionador de su propia conducta institucional (la paradoja a la que se aludió en clase).
\end{itemize}

\subsection{Registro y certificación de la clase}

Existe un \textbf{registro nacional de acciones de clase}, en el marco de la página de la Corte, que sirve para llevar un control de la acción y para especificar cómo está constituida la clase. Al registrar una acción es necesario indicar a quién se representa. La \textbf{certificación} de la acción de clase es clave; la ley de acción de clase, o certificación, es muy concisa.

\begin{example}{Cómo se define la clase}
``Estoy representando a aquellos compradores de camionetas Hilux realizadas en la planta de la Ciudad Autónoma de Buenos Aires entre los años 22 y 24'' (ejemplo hipotético del docente). Hay una acción de clase porque existe un interés homogéneo a partir del cual hay reclamos.
\end{example}

%------------------------------------------------------------------
\section{Cuadro Cornell: derechos colectivos y acciones de clase}
%------------------------------------------------------------------

\begin{longtable}{@{}p{0.26\textwidth}p{0.70\textwidth}@{}}
\toprule
\textbf{Preguntas / palabras clave} & \textbf{Notas / respuestas} \\
\midrule
\endfirsthead
\toprule
\textbf{Preguntas / palabras clave} & \textbf{Notas / respuestas} \\
\midrule
\endhead

\textbf{¿Qué es la legitimación activa y qué categorías de derechos distingue Halabi?} &
Es presupuesto de una cuestión justiciable. Halabi distingue tres categorías: derecho individual subjetivo, derechos de incidencia colectiva e intereses individuales homogéneos (causa común), estos últimos origen de las acciones de clase. \\[0.8em]

\textbf{¿Qué es una acción de clase y quién puede iniciarla?} &
Acción que inician el defensor del pueblo, asociaciones o el afectado en representación de su clase, sobre derechos con causa homogénea: divisibles en el daño, pero con origen común (ejemplo: falla en camionetas fabricadas en cierta planta y período). \\[0.8em]

\textbf{¿Para qué sirve y cómo se define la clase?} &
Sirve para disciplinar a grandes corporaciones; nunca la inicia el propio Estado contra su gestión. Se registra en el registro nacional, exige certificación y una definición precisa de quién integra la clase. \\

\midrule
\multicolumn{2}{@{}p{0.96\textwidth}@{}}{\textbf{Resumen:} La legitimación activa y las categorías de Halabi explican quién puede reclamar y por qué existen las acciones de clase, que equilibran el poder de las corporaciones frente a reclamos individuales de poco peso.} \\
\bottomrule
\end{longtable}

%------------------------------------------------------------------
\chapter{Ambiente, consumo y competencia}
%------------------------------------------------------------------

\section{Derecho ambiental}

Se deben repasar la \textbf{ley de presupuestos mínimos medioambientales} y los principios generales en materia ambiental.

La Constitución proclama un derecho general a tener y desenvolverse en un medio ambiente sano, y establece las competencias: el Congreso dicta los presupuestos mínimos (lo ha hecho mediante la ley de presupuestos mínimos), y las provincias pueden complementar la regulación, siempre en pos de maximizar los derechos relativos al medio ambiente. Es un derecho eminentemente local, con una idea de competencia concurrente entre Nación y provincias.

\begin{quote}
\textbf{Constitución Nacional, art. 41.} Todos los habitantes gozan del derecho a un ambiente sano, equilibrado, apto para el desarrollo humano y para que las actividades productivas satisfagan las necesidades presentes sin comprometer las de las generaciones futuras; y tienen el deber de preservarlo. El daño ambiental generará prioritariamente la obligación de recomponer, según lo establezca la ley. Las autoridades proveerán a la protección de este derecho, a la utilización racional de los recursos naturales, a la preservación del patrimonio natural y cultural y de la diversidad biológica, y a la información y educación ambientales. Corresponde a la Nación dictar las normas que contengan los presupuestos mínimos de protección, y a las provincias, las necesarias para complementarlas, sin que aquéllas alteren las jurisdicciones locales. Se prohíbe el ingreso al territorio nacional de residuos actual o potencialmente peligrosos, y de los radiactivos.
\end{quote}

\subsection{Evaluación de impacto ambiental}

\begin{definition}{Evaluación de impacto ambiental (EIA)}
Análisis sobre el impacto en el medio ambiente, con la cuantificación o predicción de los riesgos ambientales y las acciones de mitigación. Es la principal herramienta jurídica y científica para resolver los dilemas que se presentan entre el uso y goce de los recursos agotables del planeta y el derecho a un medio ambiente sano.
\end{definition}

\subsection{El fallo Atuel}

El fallo \textbf{Atuel} es una controversia entre dos provincias sobre el aprovechamiento de recursos hídricos compartidos. Fue incorporado como pronunciamiento nuevo, de la última camada de fallos destacables en materia ambiental, en el que la Corte reflexiona sobre muchos de los principios de la materia. \textbf{No se toma como jurisprudencia para cuestiones de hecho} (quién es actor, quién demandado, qué se resolvió), sino como complemento para ver la doctrina de la Corte en materia ambiental.

\subsection{Amparo ambiental y principio preventivo}

\begin{definition}{Amparo ambiental}
Mecanismo rápido y expedito, mediante un proceso sumarísimo, muy ligado al \textbf{principio preventivo}. En materia ambiental hay una idea de dar deferencia a las medidas que puedan considerarse adecuadas para mantener el statu quo y evitar una mayor degradación del medio ambiente. Debe demostrarse, al menos en un grado conjetural, una eventual afectación del medio ambiente.
\end{definition}

\section{Derechos del consumidor}

La \textbf{ley 24.240} es un punto muy importante. Sus ideas centrales son:
\begin{itemize}
    \item \textbf{Responsabilidad solidaria} en toda la cadena.
    \item Protección de la adquisición de un bien o servicio para aprovechamiento personal o del grupo familiar, aun cuando sea de carácter oneroso.
    \item Conceptos de \textbf{solidaridad}, \textbf{prevención} y \textbf{reparación integral}.
    \item Un plexo normativo de \textbf{orden público}, que puede llegar a prevalecer incluso por sobre la voluntad o el acuerdo de las partes.
\end{itemize}

\subsection{Arbitraje de consumo}

Entre los modos alternativos de resolución de conflictos previstos en la ley (se hace foco en el arbitraje) se encuentra el \textbf{arbitraje de consumo}, que se da entre el consumidor y la empresa.

\begin{definition}{Arbitraje de consumo}
Sistema \textbf{no judicial pero jurisdiccional}: si la decisión no se impugna, queda firme, y se impugna mediante un recurso de nulidad como cualquier laudo arbitral. Lo resuelve un tercero independiente, el árbitro del consumo, designado en el seno de la autoridad de aplicación (Subsecretaría de Defensa del Consumidor, en el ámbito de la Secretaría de Comercio). Requiere la \textbf{adhesión previa} del eventual demandado, el prestador de bienes y servicios.
\end{definition}

Por sus características, queda reservado a grandes empresas que voluntariamente se someten a él: telefonía celular, telefonía prepaga, televisión por cable y algunos servicios domiciliarios.

\begin{quote}
\textbf{Ley 24.240, art. 59 (Tribunales Arbitrales).} Contenido esencial: la autoridad de aplicación propicia la organización de tribunales arbitrales que actúan como amigables componedores o árbitros de derecho; el sometimiento de las partes al Sistema Nacional de Arbitraje de Consumo es voluntario y debe constar expresamente por escrito.
\end{quote}

\begin{note}
El recuadro anterior resume el contenido del artículo; para citar su texto literal conviene verificarlo en la versión vigente de la ley.
\end{note}

Se trata de un tema que apenas se mencionó en clase: basta entender que la ley prevé un sistema de resolución jurisdiccional, aunque no judicial, mediante la herramienta del arbitraje. Es de las pocas construcciones que se han mantenido en las reformas legales posteriores.

\section{Derecho de la competencia}

\subsection{Autoridad de aplicación y regímenes preventivo y represivo}

Se deben conocer los conceptos generales de la \textbf{ley antitrust} y las características de la autoridad de aplicación, la autoridad de la competencia: sus miembros duran cinco años en funciones (dos de ciencias económicas, dos de derecho y un tercero) y forman un tribunal que resuelve las cuestiones de ambos regímenes.
% TODO: contenido incompleto o ambiguo en las notas originales (composición de la autoridad: <<dos de ciencias económicas, dos de materia de derecho y un tercero>>)

\begin{table}[H]
\centering
\renewcommand{\arraystretch}{1.3}
\begin{tabularx}{\textwidth}{
    >{\raggedright\arraybackslash}p{0.19\textwidth}
    >{\raggedright\arraybackslash}X
    >{\raggedright\arraybackslash}X
}
\toprule
\textbf{Fase} & \textbf{Qué hace} & \textbf{Resultado posible} \\
\midrule
\textbf{Régimen preventivo} & Autorización de fusiones o adquisiciones empresariales & Declarar la autorización, condicionarla o rechazarla \\
\textbf{Fase represiva} & Aplicación de multas ante inconductas o actos que atentan contra el interés económico general & Sanciones, una vez acreditada la conducta \\
\textbf{Fase reparatoria} & Interviene cuando existe algún pronunciamiento & Reparación, según lo mencionado en clase \\
\bottomrule
\end{tabularx}
\caption{Fases del régimen de defensa de la competencia}
\end{table}

También se aclaró que no se encuentra vigente la idea del pago previo.
% TODO: contenido incompleto o ambiguo en las notas originales (alcance de la aclaración sobre el pago previo en materia de competencia)

\subsection{Posición dominante e infracciones}

\begin{definition}{Posición dominante}
Según la ley, es la situación en que un oferente de un bien o servicio no está expuesto a una real competencia. Es una buena aproximación, pero resulta un tanto tautológica; por ello conviene profundizar en la posición dominante, en el \textbf{abuso} de la posición dominante y en ejemplos claros de infracciones.
\end{definition}

Se debe repasar las distintas infracciones que la norma prevé taxativamente:
\begin{itemize}
    \item \textbf{reparto de mercados} y \textbf{cartelización}, que son las típicas;
    \item \textbf{negativas de venta};
    \item \textbf{precios predatorios};
    \item \textbf{ventas atadas};
    \item otros conceptos más sofisticados que aparecen con el crecimiento del mercado empresarial.
\end{itemize}

También interesan las posibilidades sancionatorias con que cuenta el Estado. Como guía de estudio se recomienda tomar la ley, que es ilustrativa y es un buen resumen (incluida la composición de la autoridad de aplicación), y profundizar en los puntos vistos.

\subsection{El procedimiento antitrust}

Deben comprenderse los principios básicos del debido proceso del procedimiento antitrust:
\begin{itemize}
    \item \textbf{Iniciativa}: por denuncia de terceros o de la autoridad de aplicación.
    \item \textbf{Instrucción}: se asemeja mucho a la instrucción penal de primera instancia, con todas las oportunidades para ejercer las defensas.
    \item \textbf{Imputación clara}, con los puntos de las infracciones que se intenta demostrar al sujeto investigado.
    \item Es un procedimiento entre la \textbf{autoridad de aplicación} y el \textbf{sujeto denunciado}: el denunciante puede promover la denuncia, pero luego la eventual impugnación es siempre entre la parte sancionada y la autoridad de aplicación. El denunciante ya no es parte.
\end{itemize}

%------------------------------------------------------------------
\section{Cuadro Cornell: ambiente, consumo y competencia}
%------------------------------------------------------------------

\begin{longtable}{@{}p{0.26\textwidth}p{0.70\textwidth}@{}}
\toprule
\textbf{Preguntas / palabras clave} & \textbf{Notas / respuestas} \\
\midrule
\endfirsthead
\toprule
\textbf{Preguntas / palabras clave} & \textbf{Notas / respuestas} \\
\midrule
\endhead

\textbf{¿Cómo se organiza la protección ambiental?} &
La Constitución proclama el derecho a un ambiente sano; la Nación dicta los presupuestos mínimos y las provincias los complementan para maximizar el derecho (competencia concurrente). La evaluación de impacto ambiental es la herramienta clave: analiza impactos, cuantifica riesgos y prevé mitigaciones. \\[0.8em]

\textbf{¿Qué es el amparo ambiental?} &
Proceso sumarísimo, rápido y expedito, ligado al principio preventivo, con afectación demostrada al menos conjeturalmente. El fallo Atuel se toma como complemento doctrinario, no por sus hechos. \\[0.8em]

\textbf{¿Cómo protege la ley 24.240 al consumidor?} &
Establece responsabilidad solidaria en toda la cadena, protege la adquisición de bienes y servicios para uso personal o familiar, y su normativa de orden público puede prevalecer sobre la voluntad de las partes (solidaridad, prevención, reparación integral). \\[0.8em]

\textbf{¿Qué es el arbitraje de consumo?} &
Sistema no judicial pero jurisdiccional (art. 59): árbitro independiente designado en la autoridad de aplicación, laudo que queda firme si no se impugna, y adhesión previa y voluntaria del proveedor (grandes empresas: telefonía, televisión por cable, servicios domiciliarios). \\[0.8em]

\textbf{¿Cómo funciona la defensa de la competencia?} &
La autoridad (miembros por cinco años) resuelve el régimen preventivo (autorización, condicionamiento o rechazo de fusiones y adquisiciones) y el represivo (multas por conductas que atentan contra el interés económico general). Posición dominante: oferente no expuesto a una real competencia. Infracciones típicas: reparto de mercados y cartelización, negativas de venta, precios predatorios, ventas atadas. \\[0.8em]

\textbf{¿Cómo es el procedimiento antitrust?} &
Denuncia o iniciativa de oficio, instrucción con posibilidad de defensas, imputación clara; es un procedimiento entre la autoridad y el denunciado, y el denunciante deja de ser parte. \\

\midrule
\multicolumn{2}{@{}p{0.96\textwidth}@{}}{\textbf{Resumen:} El derecho ambiental (presupuestos mínimos, evaluación de impacto y amparo preventivo), la defensa del consumidor (responsabilidad solidaria y arbitraje de consumo) y la defensa de la competencia (control preventivo y represivo del abuso de posición dominante) son aplicaciones concretas de los derechos colectivos.} \\
\bottomrule
\end{longtable}

%------------------------------------------------------------------
\chapter{Alcance y modalidad de la evaluación}
%------------------------------------------------------------------

\section{Contenidos comprendidos y excluidos}

\begin{itemize}
    \item \textbf{Comprendidos:} libertad religiosa, derecho de réplica, libertad de prensa (a nivel conceptual: doble dimensión y prohibición de censura previa), propiedad y expropiación, legitimación activa, ambiente, consumidor y competencia.
    \item \textbf{Excluidos:} los \textbf{derechos políticos}, que se vieron y repasaron pero quedan fuera de este parcial (hay tres fallos subidos al campus, pero no habrá una pregunta específica); y los temas de la primera parte ya evaluados: New York Times v. Sullivan, Real Malicia y doctrina Campillay.
\end{itemize}

\section{Modalidad y nivel de exigencia}

\begin{itemize}
    \item Examen \textbf{escrito y presencial}, con una mezcla de preguntas más puntuales y otras más para desarrollar, en la tónica del parcial anterior.
    % TODO: contenido incompleto o ambiguo en las notas originales (la fecha del examen aparece como <<jueves>> y luego como <<lunes 16>>)
    \item Se exigirá mayor \textbf{precisión} en el relato y en las respuestas: definiciones más concretas y ágiles, y un manejo de los conceptos espontáneo y oportuno.
    \item En competencia, es probable que se pida desarrollar tres o cuatro conceptos que expliquen qué es una conducta anticompetitiva, no necesariamente de forma taxativa.
    \item En arbitraje de consumo no se espera una pregunta detallada del sistema.
\end{itemize}

%------------------------------------------------------------------
\section{Cuadro Cornell: orientaciones para la evaluación}
%------------------------------------------------------------------

\begin{longtable}{@{}p{0.26\textwidth}p{0.70\textwidth}@{}}
\toprule
\textbf{Preguntas / palabras clave} & \textbf{Notas / respuestas} \\
\midrule
\endfirsthead
\toprule
\textbf{Preguntas / palabras clave} & \textbf{Notas / respuestas} \\
\midrule
\endhead

\textbf{¿Qué temas entran en el segundo parcial?} &
Libertad religiosa, réplica, libertad de prensa a nivel conceptual, propiedad, legitimación, ambiente, consumidor y competencia. \\[0.8em]

\textbf{¿Qué queda fuera?} &
Los derechos políticos y los temas del primer parcial ya evaluados (Real Malicia, Campillay). \\[0.8em]

\textbf{¿Cómo será el examen y qué se exige?} &
Escrito y presencial, con mezcla de preguntas puntuales y de desarrollo; se exige precisión conceptual, definiciones concretas y ágiles, y fundamento en la jurisprudencia vista. \\

\midrule
\multicolumn{2}{@{}p{0.96\textwidth}@{}}{\textbf{Resumen:} El examen será escrito y exigirá definiciones precisas, ágiles y fundadas en la jurisprudencia vista: Castillo, Ekmekdjian, Bourdieu y Halabi, junto con las normas de los arts. 2, 17, 41 y 43 de la Constitución y la ley 24.240.} \\
\bottomrule
\end{longtable}

\end{document}
